\documentclass[11pt]{article}
\usepackage[a4paper,margin=25mm]{geometry}
\usepackage{amsmath,amssymb,bm,booktabs}
\usepackage[T1]{fontenc}
\usepackage{hyperref}
\newcommand{\dtzero}{\partial_0}
\newcommand{\dd}{\mathrm{d}}
\newcommand{\tf}{\mathrm{TF}}

\title{A From-Scratch Characteristic Decomposition and Boundary Treatment\\
for AthenaK's Background-Subtracted Z4c System}
\author{Implementation note for \texttt{z4c\_tov\_ks}}
\date{\today}

\begin{document}
\maketitle

\begin{abstract}
This note derives the face-normal characteristic variables used by AthenaK's
residual-Z4c boundary kernel directly from the equations implemented in
\texttt{src/z4c/z4c\_calcrhs.cpp}.  No characteristic formula is imported from
a paper.  Starting from the zero-residual Fr\'echet symbol, the derivation
retains the distinct full-state and analytic-background shift advections and
obtains the exact frozen-coefficient quasilinear symbol at a finite residual.
It produces ten incoming propagating modes at each open face: four scalar,
four vector, and two tensor modes.  The note also specifies the
method-of-lines compatibility condition, field projections, finite
differences, fused GPU ownership, input validation, and the limitations of the
present \(L=1\), planar treatment.
\end{abstract}

\section{Implemented residual equations and scope}

Let a bar denote the analytic background and let
\[
  \delta u = u-\bar u
\]
denote the stored residual.  For the geometric Z4c variables AthenaK evaluates
the residual right-hand side as
\[
  {\cal R}_{\rm res}(\delta u)
  ={\cal R}(\bar u+\delta u)-{\cal R}(\bar u),
  \label{eq:residual}
\]
using identical finite-difference stencils for both terms.  Matter is included
only in the full term and is required to be vacuum or atmosphere at a physical
boundary.  The supported gauge path is the code's
\texttt{background\_adapted} path,
\begin{align}
 \dtzero\,\delta\alpha &=-L\,\delta\hat K
                  +\hbox{lower-order damping}, \label{eq:lapse}\\
 \dtzero\,\delta\beta^i &=G\,\delta\tilde\Gamma^i
                  +\hbox{lower-order damping}, \label{eq:shift}
\end{align}
where, at zero residual,
\begin{align}
 \dtzero &\equiv \partial_t-\bar\beta^i\partial_i,\\
 N&\equiv\bar\alpha,\qquad C\equiv\bar\chi,\\
 L&\equiv f_{\rm res}\,\bar\alpha
       \left(f_{\rm oplog}f_{\rm harmonicf}
             +f_{\rm harmonic}\bar\alpha\right),\\
 G&\equiv \texttt{shift\_Gamma}.
\end{align}
The present validity domain has unit lapse and shift advection, no telegraph
lapse, \texttt{shift\_H=0}, \texttt{shift\_alpha2Gamma=0}, and no
time-dependent slow-start shift driver.  The production choice has
\(\chi=\psi^{-4}\), \(L=2N\) asymptotically, and \(G=1\).

The first calculation below is the exact principal symbol of the Fr\'echet
derivative
\[
  D{\cal R}_{\rm res}(0)\,[\delta u].
  \label{eq:frechet}
\]
It supplies a simple independently checkable base point.  The implemented
kernel also uses the finite-residual extension derived next.  In both cases an
exact analytic background is invariant, and boundary data are never obtained
by extrapolating the large background fields.

\subsection{Audit of the finite-residual principal symbol}

It is important not to confuse Equation~\eqref{eq:frechet} with the
quasilinear principal symbol at an arbitrary, finite residual.  The geometric
part of AthenaK evaluates the full-state operator minus a prescribed
background source.  Its principal coefficients at finite residual are
therefore the reconstructed full lapse, conformal factor, metric, and shift.
The background-adapted gauge equations, on the other hand, deliberately
advect the residual lapse and shift with the background shift.  Consequently
the exact finite-residual symbol has two normal advection coefficients,
\[
  b_f=s_i\beta_{\rm full}^i,\qquad
  b_b=s_i\bar\beta^i,
\]
In this subsection
\[
 C=\chi_{\rm full},\qquad N=\alpha_{\rm full},
\]
while \(L\) remains the analytic-background-adapted lapse driver and \(G\)
remains the configured residual shift driver.  The symbol is not obtained by
simply replacing every barred field in the zero-residual formulas by a full
field.

This can be checked directly before computing any eigenvectors.  In a
full-conformal-metric orthonormal frame, add \(b_f\) to the diagonal entries of the
geometric scalar variables and add \(b_b\) to the \(d_\alpha\) and
\(d_{\beta s}\) entries.  The scalar characteristic polynomial becomes
\begin{align}
 \det(\lambda I-A_S^{\rm finite})
={}&\big[(\lambda-b_f)^2-CN^2\big]^2\\
 &\times\big[(\lambda-b_f)(\lambda-b_b)-CL\big]\\
 &\times\big[(\lambda-b_f)(\lambda-b_b)-\tfrac43G\big].
 \label{eq:finite-scalar-polynomial}
\end{align}
The vector polynomial is
\[
 \det(\lambda I-A_V^{\rm finite})
 =\big[(\lambda-b_f)^2-CN^2\big]
  \big[(\lambda-b_f)(\lambda-b_b)-G\big],
 \label{eq:finite-vector-polynomial}
\]
while the tensor roots remain \(b_f\pm N\sqrt C\).  Thus, for example, the
finite-residual lapse roots are
\[
 \lambda_{L,\pm}=\frac{b_f+b_b}{2}
  \mathbin{\pm}\frac12\sqrt{(b_f-b_b)^2+4CL}.
\]
These factorizations and the left eigensystem in
Section~\ref{sec:finite-rows} were derived from the two advections actually
present in \texttt{z4c\_calcrhs.cpp}.  The implementation reconstructs the
full state before the volume RHS, builds its face frame from the full
conformal metric, and evaluates \(C=\chi_{\rm full}\),
\(N=\alpha_{\rm full}\), and \(b_f\) from that state.  Only \(L\) and \(b_b\)
come from the analytic-background-adapted gauge.  It therefore applies the
exact finite-residual frozen principal projectors, rather than merely
substituting full fields into the zero-residual rows.

\section{Face frame and normal reduction}

At a boundary point choose the outward full-conformal-unit covector \(s_i\),
\[
 \tilde\gamma_{\rm full}^{ij}s_i s_j=1,
\]
and two deterministic orthonormal tangent vectors \(q_A^i\), \(A=1,2\).
All components below are in this local frozen-coefficient frame.  Thus the
full conformal metric is \(\delta_{ij}\), while the factors \(N\) and \(C\)
are retained.  Define
\[
 \partial_s=s^i\partial_i,\qquad
 b_f=s_i\beta_{\rm full}^i,\qquad b_b=s_i\bar\beta^i.
\]
Only normal derivatives are retained in the principal symbol.

Write
\[
 k=\delta\hat K,\quad
 \theta=\delta\Theta,\quad
 A_{ij}=\delta\tilde A_{ij},\quad
 \Gamma_i=\delta_{ij}\delta\tilde\Gamma^j,
\]
and introduce the normal reduction variables
\begin{align}
 d_\chi&=\partial_s\delta\chi,&
 d_\alpha&=\partial_s\delta\alpha,&
 d_{\beta i}&=\partial_s\delta\beta_i,\\
 d_{ss}&=\partial_s\left(
    \delta\tilde\gamma_{ss}
    -\frac13\delta\tilde\gamma^i{}_i\right),&
 d_{sA}&=\partial_s\delta\tilde\gamma_{sA},&
 d_{AB}^{\tf}&=\partial_s\delta\tilde\gamma_{AB}^{\tf}.
\end{align}
The algebraic trace is projected out of \(A_{ss}\) in the implementation:
\[
 A_{ss}\longrightarrow
 A_{ss}-\frac13 A^i{}_i.
\]
The trace and the remaining purely tangential zero-speed fields are not
prescribed by this normal boundary system.

\section{Principal geometric terms from the code}

\subsection{Conformal and physical Ricci terms}

The second-derivative terms assembled by
\texttt{ComputeGeometryData} give, in the one-dimensional normal reduction,
\begin{align}
 \delta\tilde R_{ss}
   &=\partial_s\Gamma_s-\frac12\partial_s^2
       \delta\tilde\gamma_{ss},\\
 \delta\tilde R_{sA}
   &=\frac12\partial_s\Gamma_A-\frac12\partial_s^2
       \delta\tilde\gamma_{sA},\\
 \delta\tilde R_{AB}
   &=-\frac12\partial_s^2\delta\tilde\gamma_{AB}.
\end{align}
For \(\chi=\psi^{-4}\), \(\phi=-\tfrac14\log\chi\).  Keeping only
second derivatives in the code's \(R^\phi_{ij}\) gives
\[
 \delta R^\phi_{ij}
 =\frac{1}{2C}\partial_i\partial_j\delta\chi
  +\frac{1}{2C}\delta_{ij}\partial_s^2\delta\chi,
\]
and hence
\[
 \delta R^\phi_{ss}=\frac1C\partial_s^2\delta\chi,\qquad
 \delta R^\phi_{AB}=\frac{1}{2C}\delta_{AB}
                    \partial_s^2\delta\chi,\qquad
 \delta R^\phi=\frac2C\partial_s^2\delta\chi.
\]
The principal Hamiltonian term used in the \(\Theta\) equation is therefore
\[
 \delta H
 =C\,\partial_s\Gamma_s+2\,\partial_s^2\delta\chi,
 \label{eq:hamiltonian-principal}
\]
after imposing the conformal determinant projection.

\subsection{Principal evolution equations}

Reading the coefficients from the residual RHS and using the Ricci terms
above yields
\begin{align}
 \dtzero k
   &=-C\,\partial_s d_\alpha,\\
 \dtzero\theta
   &=\frac12NC\,\partial_s\Gamma_s+N\,\partial_s d_\chi,\\
 \dtzero A_{ss}
   &=\frac23NC\,\partial_s\Gamma_s
     +\frac13N\,\partial_s d_\chi
     -\frac12NC\,\partial_s d_{ss}
     -\frac23C\,\partial_s d_\alpha,\\
 \dtzero\Gamma_s
   &=-\frac43N\,\partial_s k-\frac23N\,\partial_s\theta
     +\frac43\partial_s d_{\beta s},\\
 \dtzero d_\chi
   &=\frac23CN\,\partial_s k+\frac43CN\,\partial_s\theta
     -\frac23C\,\partial_s d_{\beta s},\\
 \dtzero d_{ss}
   &=-2N\,\partial_s A_{ss}+\frac43\partial_s d_{\beta s},\\
 \dtzero d_\alpha
   &=-L\,\partial_s k,\\
 \dtzero d_{\beta s}
   &=G\,\partial_s\Gamma_s.
 \label{eq:scalar-evolution}
\end{align}
For each tangent direction,
\begin{align}
 \dtzero A_{sA}
   &=\frac12NC\,\partial_s\Gamma_A
     -\frac12NC\,\partial_s d_{sA},\\
 \dtzero\Gamma_A&=\partial_s d_{\beta A},\\
 \dtzero d_{sA}&=-2N\,\partial_s A_{sA}
                       +\partial_s d_{\beta A},\\
 \dtzero d_{\beta A}&=G\,\partial_s\Gamma_A.
 \label{eq:vector-evolution}
\end{align}
For each of the two tangential trace-free polarizations,
\begin{align}
 \dtzero A_{AB}^{\tf}
   &=-\frac12NC\,\partial_s d_{AB}^{\tf},\\
 \dtzero d_{AB}^{\tf}
   &=-2N\,\partial_s A_{AB}^{\tf}.
 \label{eq:tensor-evolution}
\end{align}

\section{Block principal symbols}

Define
\[
 U_S=(k,\theta,A_{ss},\Gamma_s,d_\chi,d_{ss},
      d_\alpha,d_{\beta s})^T.
\]
Equations~\eqref{eq:scalar-evolution} have
\(\dtzero U_S=M_S\partial_s U_S\), where
{\small
\[
M_S=
\begin{pmatrix}
0&0&0&0&0&0&-C&0\\
0&0&0&NC/2&N&0&0&0\\
0&0&0&2NC/3&N/3&-NC/2&-2C/3&0\\
-4N/3&-2N/3&0&0&0&0&0&4/3\\
2CN/3&4CN/3&0&0&0&0&0&-2C/3\\
0&0&-2N&0&0&0&0&4/3\\
-L&0&0&0&0&0&0&0\\
0&0&0&G&0&0&0&0
\end{pmatrix}.
\]
}
With \(U_V=(A_{sA},\Gamma_A,d_{sA},d_{\beta A})^T\) and
\(U_T=(A_{AB}^{\tf},d_{AB}^{\tf})^T\),
\[
M_V=
\begin{pmatrix}
0&NC/2&-NC/2&0\\
0&0&0&1\\
-2N&0&0&1\\
0&G&0&0
\end{pmatrix},
\qquad
M_T=
\begin{pmatrix}
0&-NC/2\\
-2N&0
\end{pmatrix}.
\]
Direct determinants, not a literature formula, give
\begin{align}
 \det(\lambda I-M_S)
  &=(\lambda^2-CL)(\lambda^2-\tfrac43G)
    (\lambda^2-CN^2)^2,\\
 \det(\lambda I-M_V)
  &=(\lambda^2-G)(\lambda^2-CN^2),\\
 \det(\lambda I-M_T)
  &=\lambda^2-CN^2.
\end{align}
The normal characteristic speeds relative to \(\dtzero\) are consequently
\[
 c_L=\sqrt{CL},\qquad
 c_{SL}=\sqrt{\frac43G},\qquad
 c_{ST}=\sqrt G,\qquad
 c_P=N\sqrt C.
\]

\section{Zero-residual characteristic covectors}

Let \(\epsilon=\pm1\) label eigenvalue \(+\!c\) or \(-\!c\).
The following are left eigenvectors, so
\[
 w^\epsilon=\ell^\epsilon U,\qquad
 \ell^\epsilon M=\epsilon c\,\ell^\epsilon.
\]
Arbitrary nonzero rescalings of individual rows are immaterial.

\subsection{Scalar modes}

Define the cone-separation denominators
\[
 D_L=3CL-4G,\qquad D_P=3CN^2-4G.
\]
The lapse pair is
\[
 w_L^\epsilon=-\epsilon\sqrt{\frac LC}\,k+d_\alpha.
\]
The longitudinal-shift pair is
\begin{align}
 w_{SL}^{\epsilon}
 ={}&\frac{4GN}{D_L}k+\frac{2GN}{D_P}\theta
 +\epsilon\frac{2\sqrt{3G}(CN^2-G)}{D_P}\Gamma_s\\
 &+\epsilon\frac{\sqrt{3G}N^2}{D_P}d_\chi
 -\epsilon\frac{2\sqrt{3G}CN}{D_L}d_\alpha
 +d_{\beta s}.
\end{align}
Two independent light/constraint pairs are
\begin{align}
 w_{P1}^{\epsilon}
   &=\epsilon\sqrt C\,\theta+\frac C2\Gamma_s+d_\chi,\\
 w_{P2}^{\epsilon}
   &=\epsilon\left(
       \frac{4}{3\sqrt C}k+\frac{2}{3\sqrt C}\theta
       -\frac{2}{\sqrt C}A_{ss}\right)-\Gamma_s+d_{ss}.
\end{align}
\subsection{Vector and tensor modes}

For each \(A=1,2\),
\begin{align}
 w_{ST,A}^{\epsilon}
   &=\epsilon\sqrt G\,\Gamma_A+d_{\beta A},\\
 w_{P,A}^{\epsilon}
   &=-\epsilon\frac{2}{\sqrt C}A_{sA}-\Gamma_A+d_{sA}.
\end{align}
For each of the \(+\) and \(\times\) TT polarizations,
\[
 w_{TT}^{\epsilon}
   =-\epsilon\frac{2}{\sqrt C}A_{AB}^{\tf}+d_{AB}^{\tf}.
\]
These comprise four scalar, four vector, and two tensor incoming members.

\section{Exact finite-residual covectors}
\label{sec:finite-rows}

The finite-residual rows follow directly by solving the left nullspaces of
the matrices audited in
Equations~\eqref{eq:finite-scalar-polynomial} and
\eqref{eq:finite-vector-polynomial}.  For any positive driver \(Q\), define
\begin{align}
 \lambda_Q^\epsilon
 &=\frac{b_f+b_b}{2}
   +\frac{\epsilon}{2}\sqrt{(b_f-b_b)^2+4Q},\\
 \mu_Q^\epsilon&=\lambda_Q^\epsilon-b_f,\qquad
 \delta_Q^\epsilon=\lambda_Q^\epsilon-b_b.
\end{align}
Thus \(\mu_Q^\epsilon\delta_Q^\epsilon=Q\).  The two physical roots are
\[
 \lambda_P^\epsilon=b_f+\epsilon N\sqrt C.
\]

\subsection{Finite scalar rows}

The lapse row is
\[
 w_L^\epsilon=-\frac{\delta_{CL}^\epsilon}{C}k+d_\alpha.
\]
For the longitudinal shift put \(Q=4G/3\),
\(\mu=\mu_Q^\epsilon\), and \(\delta=\delta_Q^\epsilon\).  A row normalized
to unit \(d_{\beta s}\) coefficient is
\begin{align}
 w_{SL}^\epsilon={}&
 \frac{N\delta^2}{CL-\mu\delta}\,k
 +\frac{N\mu\delta}{2(CN^2-\mu^2)}\,\theta\\
 &+\frac{\delta(4CN^2-3\mu^2)}
             {4(CN^2-\mu^2)}\,\Gamma_s
 +\frac{N^2\delta}{2(CN^2-\mu^2)}\,d_\chi\\
 &-\frac{CN\delta}{CL-\mu\delta}\,d_\alpha+d_{\beta s}.
\end{align}
The GPU implementation multiplies this row by
\((CL-\mu\delta)(CN^2-\mu^2)\).  This leaves the eigenspace and boundary
condition unchanged while avoiding large intermediate coefficients near a
supported but nondegenerate cone separation.
The two light rows are unchanged except that their coordinate eigenvalues
are \(\lambda_P^\epsilon\):
\begin{align}
 w_{P1}^{\epsilon}
   &=\epsilon\sqrt C\,\theta+\frac C2\Gamma_s+d_\chi,\\
 w_{P2}^{\epsilon}
   &=\epsilon\left(
       \frac{4}{3\sqrt C}k+\frac{2}{3\sqrt C}\theta
       -\frac{2}{\sqrt C}A_{ss}\right)-\Gamma_s+d_{ss}.
\end{align}
Setting \(b_f=b_b\) reduces these rows algebraically to the zero-residual
rows above, up to arbitrary nonzero normalization.

\subsection{Finite vector and tensor rows}

For each tangent direction the transverse-shift roots use \(Q=G\):
\[
 w_{ST,A}^{\epsilon}
 =\delta_G^\epsilon\Gamma_A+d_{\beta A}.
\]
The constraint and radiative rows retain their zero-residual form,
\begin{align}
 w_{P,A}^{\epsilon}
   &=-\epsilon\frac{2}{\sqrt C}A_{sA}-\Gamma_A+d_{sA},\\
 w_{TT}^{\epsilon}
   &=-\epsilon\frac{2}{\sqrt C}A_{AB}^{\tf}+d_{AB}^{\tf},
\end{align}
with coordinate eigenvalue \(\lambda_P^\epsilon\).

\section{Constraint and radiative meaning of the light-speed rows}

The light-speed rows above can be related directly to the principal Z4
constraint subsystem; this is a useful check that they are not merely
light-speed eigenvectors of the main evolution system.  In the frozen
full-conformal orthonormal frame, the conformal connection constraint is
\[
 Z_s=\frac12(\Gamma_s-d_{ss}),\qquad
 Z_A=\frac12(\Gamma_A-d_{sA}).
\]
The normal principal parts of the Hamiltonian and momentum constraints are
\begin{align}
 H&=C\,\partial_s\Gamma_s+2\,\partial_s d_\chi,\\
 M_s&=\partial_s A_{ss}
      -\frac23\partial_s(k+2\theta),\\
 M_A&=\partial_s A_{sA}.
\end{align}
Taking a normal derivative of the incoming (\(\epsilon=+1\)) light-speed
rows gives the identities
\begin{align}
 2\,\partial_s w_{P,1}^{+}
   &=H+2\sqrt C\,\partial_s\theta,\\
 \partial_s w_{P,2}^{+}
   &=-\frac{2}{\sqrt C}\left(M_s+\partial_s\theta\right)
     -2\,\partial_s Z_s,\\
 \partial_s w_{P,A}^{+}
   &=-\frac{2}{\sqrt C}M_A-2\,\partial_s Z_A.
\end{align}
Thus the two scalar and two vector rows are spatial primitives of the four
incoming principal constraint characteristics, up to fixed nonzero
normalizations.  Prescribing them supplies the constraint sector and does
not consume an outgoing main-system characteristic.

For a tangential trace-free metric perturbation \(h_{AB}^{\tf}\), freeze the
shift for this local check and use
\[
 A_{AB}^{\tf}=-\frac{1}{2N}\partial_t h_{AB}^{\tf}
 \quad\hbox{and}\quad
 d_{AB}^{\tf}=\partial_s h_{AB}^{\tf}.
\]
The incoming radiative row is therefore
\[
 w_{TT}^{+}
 =\partial_s h_{AB}^{\tf}
  +\frac{1}{N\sqrt C}\partial_t h_{AB}^{\tf}.
\]
It vanishes for a purely outgoing normal wave.  The incoming Weyl scalar
\(\Psi_0^{\rm res}\) has, at principal order, one further outgoing-null
derivative of this row (plus tangential and lower-order curvature terms).
Consequently homogeneous \(w_{TT}^{+}=0\) is the \(L=1\) maximally
dissipative condition compatible with zero residual \(\Psi_0\).  It is
stronger than merely requiring a constant incoming waveform, but it is not
the higher-order auxiliary-field absorption hierarchy.

\section{Which sign is incoming?}

For an eigenmode of the complete finite-residual normal symbol with eigenvalue
\(\lambda\), the outward coordinate transport velocity is
\[
 v_{\rm out}=-\lambda.
\]
A negative \(v_{\rm out}\) moves from the boundary into the domain.  Therefore
the member is incoming when
\[
 \lambda>0.
\]
The code evaluates all four root pairs and requires
\(\lambda^+>0\), \(\lambda^-<0\) for each.  Thus precisely every
\(\epsilon=+1\) member is incoming: ten modes per open face.  It fails rather
than silently changing the mode count.

\section{Method-of-lines boundary compatibility}

For a perturbation whose residual has compact support away from the boundary,
the desired homogeneous continuum datum is
\[
 w^+_{\rm res}=0.
\]
The method-of-lines compatibility condition is its time derivative, not a
cell-scale damping equation:
\[
 \partial_t w^+_{\rm res}=0.
 \label{eq:target}
\]
If the initial data satisfy the homogeneous datum, this preserves it.  More
generally it preserves the incoming characteristic datum already present at
the start of the evolution.  The distinction matters when the residual
contains a stationary Coulomb or matter near field: such a solution can have
nonzero incoming and outgoing \emph{principal} amplitudes whose fluxes
balance, even though it contains no incoming dynamical wave.  Forcing that
amplitude itself to zero would change the stationary solution and launch a
boundary transient.
At the first active boundary cell the kernel forms \(w^+\) solely from stored
residual fields and their inward, second-order normal derivatives.  It also
forms the rate already supplied by the complete volume RHS and any
problem-generator source:
\[
 \dot w^+_{\rm vol}
 =\ell_p^+\,p_{\rm rhs}
  +\ell_d^+\,\partial_s q_{\rm rhs}.
 \label{eq:volume-rate}
\]
Here \(p=(k,\theta,A,\Gamma)\) and
\(q=(\chi,\tilde\gamma,\alpha,\beta)\).  Because
Equation~\eqref{eq:volume-rate} uses the already assembled residual RHS, it
retains the implemented nonlinear and lower-order constraint damping,
background subtraction, gauge damping, Kreiss--Oliger/source contributions,
and radiative terms; the boundary kernel does not reimplement them.
The left rows, frame, and characteristic speeds are held fixed while this
method-of-lines rate is formed.  This is the standard frozen-local
quasilinear boundary operator: its lower-order data are the complete
nonlinear volume/source rate, not a linearized or background-extrapolated
surrogate.  A time derivative of the local eigensystem is intentionally not
part of the frozen operator.

The implemented \(L=1\) planar target is therefore
\[
 \dot w^+_{\rm target}=0.
\]
It supplies no time-dependent incoming characteristic data.  In particular,
the implementation does \emph{not} use the rejected grid-scale relaxation
\(\dot w^+=-\lambda^+w^+/\Delta_s\).  That relaxation gives the correct fixed
point for compact perturbations but is inconsistent with a stationary
nonzero residual near field and was observed to inject a returning TDE
transient.  The adopted equation is a method-of-lines/Bj{\o}rhus
compatibility condition, not a hard overwrite of a ghost cell or active
field.

\subsection{Runtime source selection and validated provenance}

The runtime selector \texttt{boundary\_rhs} now defaults to
\texttt{characteristic\_cpbc}.  Within that boundary operator,
\texttt{characteristic\_bc\_source} selects the compatibility target:
\texttt{zero\_rate} is the default and imposes
\(\dot w^+_{\rm target}=0\), while
\texttt{tangential\_principal} selects the separate experimental target
formed from the frozen tangential and mixed principal terms.  The generic
\texttt{boundary\_rhs=characteristic\_cpbc} path therefore retains the
validated zero-rate behavior unless the experimental source is requested
explicitly.  The tangential-principal source has not been validated in the
TDE evolution and must not be identified with the result reported for the
zero-rate source.

The exact recovered \texttt{src/z4c/z4c\_Sbc.cpp} snapshot corresponding to
the validated zero-rate executable has SHA-256
\begin{center}
\texttt{5abb0ef3c22c8f989b5916d277676b0a}\\
\texttt{2a7ef993540edfe917c3e74d1b76a2c8}
\end{center}
and Git blob \texttt{437fe5be205f8a2130f85f68ec75f3500d50cb1b}.
The archived executable used for the validation has SHA-256
\begin{center}
\texttt{cd7cefef042ff075e85688dd3ec08dca}\\
\texttt{243df479d50b91e0439dd8a20da46478}.
\end{center}
These identifiers record the already tested implementation; introducing the
runtime split does not constitute a new validation of either source mode.

To preserve the code's second-order evolution structure, corrections are
applied only to the momentum-like RHS variables:
\[
 (\hat K,\Theta,\tilde A_{ij},\tilde\Gamma^i)_{\rm rhs}.
\]
The RHS values for \(\chi,\tilde\gamma_{ij},\alpha,\beta^i\), all tangential
zero-speed fields, and all auxiliary \(B_i\) remain those of the volume
operator.  In the scalar block the four equations
\[
 L_p^+\,\delta p_{\rm rhs}
 =\dot{\bm w}_{\rm target}^+-\dot{\bm w}_{\rm vol}^+
\]
are solved by the closed sparse triangular map implied by the displayed rows;
no augmented matrix or runtime eigensolve is allocated in a GPU thread.  In
each vector block,
\begin{align}
 \delta\Gamma_{A,\rm rhs}
   &=\frac{r_{ST,A}}{\delta_G^+},\\
 \delta A_{sA,\rm rhs}
   &=-\frac{\sqrt C}{2}
       \left(r_{P,A}+\delta\Gamma_{A,\rm rhs}\right),
\end{align}
and in each tensor block
\[
 \delta A_{AB,\rm rhs}^{\tf}
   =-\frac{\sqrt C}{2}r_{TT}.
\]
Here each \(r\) is target rate minus volume rate.

This prescription supplies data only for the incoming members.  ``Outgoing
unchanged'' means that their state and volume RHS are not independently
prescribed at the boundary.  Because a second-order-in-space system is updated
through its momentum variables, an \(A\)- or \(\Gamma\)-RHS correction appears
algebraically in both signs if one differentiates both characteristic
amplitudes in time.  It would be incorrect to claim that both outgoing
\emph{characteristic rates} and all configuration-variable RHSs can
simultaneously remain unchanged.  The adopted compatibility form is the
standard non-overdetermining choice for the second-order system.

\section{Cartesian reconstruction and GPU kernel}

The scalar, vector, and tensor corrections are reconstructed without a
temporary array.  For example, the trace-free curvature correction is
\begin{align}
 \delta A_{ij}^{\rm corr}
={}&\delta A_{ss}
 \left(s_i s_j-\frac12q_{1i}q_{1j}-\frac12q_{2i}q_{2j}\right)\\
&+\sum_A\delta A_{sA}(s_iq_{Aj}+q_{Ai}s_j)\\
&+\delta A_+
  (q_{1i}q_{1j}-q_{2i}q_{2j})
 +\delta A_\times(q_{1i}q_{2j}+q_{2i}q_{1j}).
\end{align}
The analogous vector reconstruction is
\[
 \delta\Gamma^i_{\rm corr}
 =\delta\Gamma_s s^i+\delta\Gamma_1q_1^i+\delta\Gamma_2q_2^i.
\]

There are three disjoint orientation kernels and at most three GPU launches
per boundary application.  Each launch handles both signs of that coordinate
face and all characteristic sectors in one thread.  Before the launches, a
host scan of the six boundary flags builds three persistent compact lists
containing only local MeshBlocks that touch a physical face in the
corresponding orientation.  The lists are copied to the device only when AMR
or load balancing changes their contents; an orientation with an empty list
is not launched.  Thus interior MeshBlocks incur no CPBC device work and no
runtime field-sized temporary is allocated.
The \(x^1\) launch owns every cell incident on an \(x^1\) physical face; the
\(x^2\) launch skips those cells, and the \(x^3\) launch skips cells already
incident on \(x^1\) or \(x^2\).  At an edge or corner the owning thread
constructs one normalized composite covector from all incident outward face
covectors.  Consequently no two kernels write the same cell even though an
edge/corner MeshBlock may occur in more than one compact list.

No analytic-background ghost extrapolation enters a characteristic
amplitude.  The local frame and geometric coefficients are read from the
reconstructed full state at the active cell; the analytic background supplies
only the gauge driver and background-shift advection identified explicitly
above.  Residual ghost extrapolation remains enabled for the centered volume
and KO stencils.
The physical ghost fill is distinct from the active-cell characteristic
correction.  For the default linear extrapolation, an outer \(x^1\) ghost
layer is
\[
 u_{i_e+\delta}=u_{i_e}
   +\delta\left(u_{i_e}-u_{i_e-1}\right),
 \qquad \delta=1,\ldots,n_g ,
\]
with analogous quadratic and cubic formulas when
\texttt{extrap\_order} is three or four.  AthenaK performs the physical
face sweeps in the fixed order \(x^1,x^2,x^3\), each over the full transverse
range including ghost indices.  Edge and corner ghosts therefore receive the
ordered tensor-product polynomial extension: an \(x^1x^2\) edge is extended
first in \(x^1\) and then in \(x^2\), while an \(x^1x^2x^3\) corner is
extended successively in all three directions.  A direction that is an
internal MeshBlock interface is instead supplied by neighbor exchange or AMR
prolongation.  The SYCL queue is in order, so these face sweeps do not race.

This separation also exposes a discrete-accuracy caveat.  Linear
extrapolation gives an \(O(h^2)\) ghost value, but inserting that value in a
centered first derivative at the last active cell can leave an \(O(h)\)
boundary derivative.  The characteristic correction removes the resulting
incoming component from the corrected momentum-like rates, but it deliberately
does not replace the outgoing and tangential volume rates.  Their local
truncation error can therefore generate a reflected pulse even when the
reported enforcement residual is at roundoff.  This construction is not
claimed to be an SBP--SAT energy estimate with a proven high-order boundary
closure.  The interior reflected-characteristic test measures that effect
independently of the boundary enforcement diagnostic; quadratic/cubic
extrapolation or a compatible one-sided/SBP derivative closure must be used if
the requested convergence and reflection gates fail.

The composite-normal treatment at an active edge or corner is a deterministic
discrete closure.  It avoids applying incompatible face projectors
sequentially, but it is not asserted to be the exact continuum
multidimensional CPBC at a measure-zero nonsmooth boundary point.  Oblique
edge/corner pulses and repeated-run checksums test this closure directly.

\section{Finite differences, checks, and reproducibility}

For an incident coordinate direction with side
\(\sigma_i=\pm1\), the active-cell derivative is
\[
 \partial_i f
 =\frac{\sigma_i}{2\Delta x_i}
  \left(3f_0-4f_{-\sigma_i}+f_{-2\sigma_i}\right).
\]
Nonincident components of a non-axis-aligned metric normal use the centered
second-order derivative.  The same rule is applied to the residual fields and
to their complete RHS.

Startup and first-use checks require finite positive \(N,C,L,G\), a
positive-definite full conformal metric (tested by all three leading principal
minors), real separated roots, exactly one incoming member of every pair, an
enrolled analytic-background callback, and the supported gauge parameters.
Every boundary component of \(E,S_i,S_{ij}\) must be finite and below the
configured atmosphere threshold.  Degenerate finite-residual lapse/shift or
light/shift separations are rejected.  Diagnostics report maxima of incoming
gauge, constraint, and radiative amplitudes, compatibility error, correction
size, shift-to-cone ratio, invalid cells, the maximum compact boundary-block
counts by orientation, fused-kernel wall time, the measured full residual-Z4c
RHS wall time, and their ratio.  Timing samples fence the ordered device queue
only when diagnostics are requested; the normal production path introduces
no timing synchronization.

The executable symbolic check is
\[
\texttt{analysis/z4c\_characteristic/derive\_residual\_characteristics.py}.
\]
It constructs both the zero-residual and hybrid finite-residual
\(M_S,M_V,M_T\) blocks from the equations displayed here and verifies:
\begin{enumerate}
 \item every closed-form left eigenpair;
 \item the three factored characteristic polynomials;
 \item \(LR-I=0\) in exact arithmetic;
 \item idempotence of the incoming projector;
 \item annihilation of that projector by every outgoing left eigenvector;
 \item all right eigenpairs; and
 \item invertibility of the four-variable scalar RHS compatibility map.
\end{enumerate}
The symbolic check proves the finite-residual characteristic polynomials and
left eigenpairs.  The dependency-light numerical companion additionally
checks \(LR-I\), projector idempotence, outgoing preservation, and right
eigenpairs below \(10^{-12}\) for randomized distinct \(b_f,b_b\).  The
samples include the production asymptotic choice \(C=N=G=1,\ L=2\), plus
actual Schwarzschild and spinning Kerr--Schild metrics and composite corner
normals.  Spin does not add a symbol parameter after transforming the
arbitrary positive-definite local conformal metric to its orthonormal face
frame.

\section{Limitations}

This is complete incoming-sector coverage for the implemented frozen residual
system, but it is an \(L=1\) planar condition.  It is not the higher-order
auxiliary-field absorption hierarchy sometimes also called a ``full''
radiative CPBC.  It assumes a stationary analytic background at the physical
boundary.  The principal projectors are exact at finite residual, but their
coefficients and frame are frozen during each local compatibility operation;
time derivatives of the projectors themselves are not added as separate
lower-order data.  The already assembled full-minus-background volume RHS
does retain the implemented nonlinear geometry, constraint damping, gauge
damping, source, and KO terms.  Coincident gauge/light cones are
mathematically diagonalizable in some parameter limits, but the displayed
closed form is singular there, so this implementation rejects those limits
instead of using an unverified limiting basis.  Matter must not reach the
boundary.

For initially homogeneous incoming data, the condition preserves zero
incoming modes of the residual system.  For a nonzero stationary residual
near field, it preserves the frozen incoming datum rather than erasing the
near field; consequently it should be described as ``no incoming dynamical
data'', not as a pointwise guarantee that every incoming principal amplitude
vanishes.  It does not make a finite outer boundary causally irrelevant:
discretization error, tangential derivatives omitted by the frozen normal
symbol, AMR interfaces, and lower-order coefficient/frame variation can still
reflect.  Those effects are measured by the plane-pulse, exact-background,
small-boundary TDE, and production smoke tests rather than assumed away.

\end{document}
